\documentclass[12pt]{article}
\usepackage[margin=1in]{geometry}
\usepackage[utf8]{inputenc}
\usepackage[english]{babel}
\usepackage{setspace}
\usepackage{url}
\usepackage{footmisc}
\usepackage{fancyhdr}

\doublespacing

\pagestyle{fancy}
\fancyhf{}
\fancyhead[R]{Vaught \thepage}

\begin{document}

\begin{center}
\Large \textbf{Classical vs. Modern Hippocratic Oath: \\
A Defense of Tradition in Contemporary Medicine}
\end{center}

\bigskip

The \textit{Hippocratic Oath}, attributed to Hippocrates and traditionally recited by new physicians for over two millennia,\footnote{Raphael Hulkower, 
``The History of the Hippocratic Oath: Outdated, Inauthentic, and Yet Still Relevant,'' 
\textit{Albert Einstein College of Medicine}, Bronx, New York 10461, 
\url{https://einsteinmed.edu/uploadedFiles/EJBM/page41_page44.pdf} (accessed January 20, 2025).} 
has been widely revised to address modern ethical dilemmas. 
In 1964, Dr.\ Louis Lasagna composed a more contemporary oath\footnote{Louis Lasagna, 
``Hippocratic Oath,'' (Boston: Tufts University School of Medicine, 1964), 
accessed January 20, 2025, \url{https://www.pbs.org/wgbh/nova/article/hippocratic-oath-today/}.} 
that explicitly emphasizes patient autonomy, social justice, and scientific advancement. 
Many argue that this modern version better suits a world marked by technological innovation, complex healthcare systems, 
and evolving moral philosophies. 
Yet upon careful comparison, the \textit{Classical Hippocratic Oath}, 
translated by Ludwig Edelstein,\footnote{Ludwig Edelstein, translator, \textit{Hippocratic Oath} 
(Baltimore: Johns Hopkins University Press, 1943), accessed January 20, 2025, 
\url{https://www.pbs.org/wgbh/nova/article/hippocratic-oath-today/}.} 
remains remarkably relevant: its call for harm avoidance, humility, and enduring moral responsibility aligns well with scientific norms 
and continues to guide physicians through contemporary ethical challenges. 
Because of its timeless principles and implicit adaptability, 
the classical oath proves no less compatible with modern medicine than its more recent counterpart.

\bigskip
\textbf{Enduring Values in the Classical Oath}

Detractors of the classical oath often point to its invocation of Greek gods (Apollo, Asclepius, Hygieia, and Panacea) 
as an example of antiquated religiosity. 
Yet this invocation can be read as a symbolic reminder of the gravity of the physician’s calling rather than a literal appeal to divine authority. 
Invoking higher powers highlights humility, underscoring that the physician’s expertise must be wielded with caution and reverence. 
Likewise, the much-discussed prohibition on administering a ``deadly drug''\footnote{Edelstein, \textit{Hippocratic Oath}.} 
does not automatically preclude nuanced discussions of end-of-life care; 
historical analyses suggest the oath aimed to prevent doctors from being coerced into assassinations or using their skills to harm for personal gain.\footnote{Baker Institute, 
``Deciphering the Hippocratic Oath in the 21st Century,'' accessed January 20, 2025, 
\url{https://www.bakerinstitute.org/research/why-hippocratic-oath-still-relevant}.} 
Rather than hampering compassion, the classical text firmly establishes a baseline of ``do no harm,'' 
which remains a bedrock ethic in medical practice.

Critics also note that the classical oath appears to forbid surgery, referencing the line about deferring operations to specialists.\footnote{Edelstein, \textit{Hippocratic Oath}.} 
However, scholars argue that this clause should be interpreted in light of ancient professional specializations. 
In reality, Greek physicians did perform invasive procedures; 
the oath’s wording likely recognized that different practitioners possessed different areas of expertise.\footnote{Wasim Shah Mushtque Shah et al., 
``History of Surgery in Ancient, Middle and Modern Era -- A Review Article,'' 
\textit{International Journal of Creative Research Thoughts (IJCRT)} 9, no. 11 (2021): 254--258, 
\url{https://ijcrt.org/papers/IJCRT2111305.pdf} (accessed January 20, 2025).} 
This emphasis on teamwork and deference to qualified colleagues actually resonates with the collaborative nature of modern healthcare, 
where radiologists, surgeons, and oncologists combine their specialties to ensure the best patient outcomes.

\bigskip
\textbf{Modern Oath: Advantages and Counterpoints}

Proponents of the modern oath emphasize its explicit discussion of patient autonomy, social context, 
and the evolving role of technology in healthcare. 
Lasagna’s text calls upon physicians to “remember that [they] do not treat a fever chart, or a cancerous growth, but a sick human being,” 
thus underscoring individualized care. 
It also highlights the need to remain open to new scientific knowledge and to share that knowledge for the betterment of patient care. 
On the surface, these provisions appear to address issues—such as informed consent and multidisciplinary research—that the classical version never explicitly names.

Nevertheless, the core values of the modern oath can be found implicitly in the Hippocratic tradition. 
The classical oath’s injunction to practice to the best of one’s ``ability and judgment'' places a responsibility on each physician 
to keep pace with medical advances, ensuring they offer no unsafe or unproven treatment. 
Similarly, its stance against injustice and exploitation translates well into modern ethics, 
where physicians are expected to prioritize patient welfare over profit and resist any undue external influences. 
Even though Lasagna’s oath provides more direct language about autonomy and public health, 
the classical oath’s general principles readily encompass these aims—particularly when supplemented by current legal and professional guidelines.

\bigskip
\textbf{Compatibility with Science and Mertonian Norms}

A crucial question is which version better aligns with modern science. 
Sociologist Robert K.\ Merton described four key norms for scientific communities: universalism, communalism, disinterestedness, 
and organized skepticism.\footnote{Robert K.\ Merton, ``A Note on Science and Democracy,'' 
\textit{Journal of Legal and Political Sociology} 1 (1942): 115--126; 
see also ``Basic Concepts: The Norms of Science,'' \textit{Adventures in Ethics and Science}, 
ScienceBlogs, January 29, 2008, \url{https://scienceblogs.com/ethicsandscience/2008/01/29/basic-concepts-the-norms-of-sc}.} 
The classical Hippocratic Oath arguably supports each of these ideals:

\begin{itemize}
    \item \textit{Universalism:} By requiring fairness and avoidance of injustice, the classical text helps prevent discrimination in care.
    \item \textit{Communalism:} The oath’s tradition of mentoring younger physicians encourages sharing knowledge for the collective good.
    \item \textit{Disinterestedness:} Physicians swear to prioritize patient welfare above personal or financial gain, mitigating conflicts of interest.
    \item \textit{Organized Skepticism:} The admonition to practice only within one’s competence and to respect specialization 
    aligns with rigorous scientific inquiry and peer review.
\end{itemize}

While the modern oath uses more contemporary language to address such themes, 
the essence of scientific integrity and ethical responsibility is already deeply rooted in the classical oath. 
Physicians guided by Hippocratic principles are, by definition, called to remain honest, collaborative, 
and dispassionate about external pressures—all values that stand at the heart of good science.

\bigskip
\textbf{Conclusion}

Ultimately, both the classical and modern oaths aim to guide physicians toward ethical and effective practice. 
Admittedly, the modern oath explicitly names autonomy, technology, and social responsibility, 
but the \textit{Classical Hippocratic Oath} still serves as a robust moral framework: 
it underscores the primacy of patient welfare, discourages harm, and promotes humility before medical complexity. 
Examined through the lens of contemporary science and Mertonian norms, its timeless principles remain surprisingly adaptable. 
Far from being an outdated relic, the classical oath can continue to inspire physicians to serve with integrity, 
collaboration, and a steadfast commitment to alleviating human suffering.

\bigskip
\begin{center}
\textbf{References}
\end{center}

\begin{itemize}
    \setlength\itemsep{0.5em}
    \item Baker Institute. ``Deciphering the Hippocratic Oath in the 21st Century.'' Accessed January 20, 2025. 
    \url{https://www.bakerinstitute.org/research/why-hippocratic-oath-still-relevant}

    \item Edelstein, Ludwig, translator. \textit{Hippocratic Oath.} Baltimore: Johns Hopkins University Press, 1943. 
    Accessed January 20, 2025. 
    \url{https://www.pbs.org/wgbh/nova/article/hippocratic-oath-today/}

    \item Hulkower, Raphael. ``The History of the Hippocratic Oath: Outdated, Inauthentic, and Yet Still Relevant.'' 
    \textit{Albert Einstein College of Medicine}, Bronx, New York 10461. 
    \url{https://einsteinmed.edu/uploadedFiles/EJBM/page41_page44.pdf} 

    \item Lasagna, Louis. ``Hippocratic Oath.'' Boston: Tufts University School of Medicine, 1964. 
    Accessed January 20, 2025. 
    \url{https://www.pbs.org/wgbh/nova/article/hippocratic-oath-today/}

    \item Merton, Robert K. ``A Note on Science and Democracy.'' \textit{Journal of Legal and Political Sociology} 1 (1942): 115--126.

    \item Mushtque Shah, Wasim Shah, et al. ``History of Surgery in Ancient, Middle and Modern Era -- A Review Article.'' 
    \textit{International Journal of Creative Research Thoughts (IJCRT)} 9, no. 11 (2021): 254--258. 
    \url{https://ijcrt.org/papers/IJCRT2111305.pdf}

    \item ScienceBlogs. ``Basic Concepts: The Norms of Science.'' \textit{Adventures in Ethics and Science}, January 29, 2008. 
    \url{https://scienceblogs.com/ethicsandscience/2008/01/29/basic-concepts-the-norms-of-sc}

\end{itemize}

\end{document}
